L'alcalde d'un poble de Catalunya encarrega a l'arquitecte municipal el disseny d'un parc
infantil que es construirà en un terreny públic. Per a complir la normativa vigent, en el parc
hi ha d'haver dos espais ben delimitats: un per a una boca de reg ---el qual, segons
l'arquitecte, ha de tenir forma circular---, i un altre per a una caseta on es guardin les
eines de manteniment ---el qual ha de tenir forma quadrada. Per motius estètics, l'arquitecte
vol delimitar cadascun d'aquests dos espais amb una barana de forja.

\begin{apartats}

\apartat{2,5}
Sabent que les dues baranes mesuren exactament 10\,m de longitud en total, quina mida ha de
tenir la barana de cada espai per tal que la suma de les superfícies dels dos espais sigui la
més petita possible? Quina és aquesta superfície mínima?

\begin{solucio}
Anomenem $x$ la longitud en metres del primer tros de barana, amb el qual es farà el tancat
circular; de l'altre tros, de $10-x$ metres, se'n farà el tancat quadrat. Com que $x$ ha de
ser el perímetre d'una circumferència, el seu radi serà $r=\frac{x}{2\pi}$, i l'àrea de
l'espai circular valdrà $A_1=\pi\left(\frac{x}{2\pi}\right)^2$. D'altra banda, com que
$10-x$ és el perímetre d'un quadrat, el seu costat serà $c=\frac{10-x}{4}$, i l'àrea de
l'espai quadrat valdrà $A_2=\left(\frac{10-x}{4}\right)^2$. La funció a minimitzar és
\[
A(x)=\pi\left(\frac{x}{2\pi}\right)^2+\left(\frac{10-x}{4}\right)^2
=\frac{x^2}{4\pi}+\frac{100-20x+x^2}{16}
=\frac{(4+\pi)x^2-20\pi x+100\pi}{16\pi}.
\]
Derivem i busquem els punts crítics:
\[
A'(x)=\frac{(8+2\pi)x-20\pi}{16\pi}=0\ \Longrightarrow\ x=\frac{20\pi}{8+2\pi}\simeq4{,}4\ \text{m}.
\]
Per assegurar que es tracta d'un mínim, la segona derivada és positiva:
$A''(x)=\dfrac{8+2\pi}{16\pi}>0$.\\
Així doncs, per tal de minimitzar la superfície total, l'arquitecte haurà de fer un
\textbf{tancat circular amb 4,4\,m} de barana i un \textbf{tancat quadrat amb els 5,6\,m}
restants. La superfície total dels dos tancats serà
\[
A(4{,}4)=\frac{(4+\pi)\cdot4{,}4^2-20\pi\cdot4{,}4+100\pi}{16\pi}\simeq3{,}5\ \text{m}^2.
\]
\textit{Pauta oficial:} 0,25 per plantejar el perímetre de la circumferència; 0,25 pel
perímetre del quadrat; 0,5 per l'expressió de l'àrea total en funció d'una sola variable;
0,5 per la derivada i l'obtenció del punt crític; 0,5 per comprovar que es tracta d'un mínim,
i 0,5 per la resposta final, incloent-hi les dues longituds i el valor de l'àrea mínima.
\end{solucio}

\end{apartats}
